\documentclass[12pt]{article}
\usepackage[utf8]{inputenc}
\usepackage[T1]{fontenc}
\usepackage{hyperref}
\title{Pamona Language Resources for K-12 Curriculum Development}
\author{Compiled by Claude Sonnet 5.5 (claude-sonnet-5-5) for Andr\'as Kornai}
\date{}
\begin{document}
\maketitle

\section{Introduction}
Pamona (also Poso or Bare'e; ISO 639-3 pmf) is an Austronesian language of Central and South Sulawesi, Indonesia, in the northern Kaili--Pamona group. The Poso dialect is the prestige variety. Estimates range from about 130,000 to 140,000 speakers, from older surveys.

\section{Primary Educational Resources}
No primers, textbooks or school programmes were found in my searches.

\section{Linguistic and Reference Materials}
\begin{itemize}
\item Pamona--Indonesian dictionary compiled by Dj.\ Tiladuru, edited by Albert Schrauwers, updating N.\ Adriani's Bare'e--Dutch dictionary to modern orthography (LOBO: Annals of Sulawesi Research, vol.\ 1; open licence; uneven draft): \url{https://lobo.journals.yorku.ca/index.php/default/article/download/4/4}
\item Wikipedia overview with phonology: \url{https://en.wikipedia.org/wiki/Pamona_language}
\item Kaili--Pamona languages overview: \url{https://en.wikipedia.org/wiki/Kaili%E2%80%93Pamona_languages}
\item Omniglot page: \url{https://omniglot.com/writing/pamona.htm}
\item An Atlantis Press paper that appeared in results (content not examined): \url{https://www.atlantis-press.com/article/125999904.pdf}
\end{itemize}

\section{Religious and Community Resources}
None found.

\section{Digital and Audio Resources}
Glosbe has community-built translations (quality unverified): \url{https://glosbe.com/pmf/mxv}.

\section{Curriculum Development Context}
The one substantial resource found is the dictionary. No dedicated modern grammar was found in these searches. Adriani's older work is the historical source, but I did not confirm its availability. Curriculum work would need local partners in Poso and would likely build on the dictionary.
\end{document}